\chapter{Relative Homology and Exact Sequences}
\label{ch:viii19-relative-homology-exact-sequences}

Relative homology measures what remains in a space after a chosen subspace is
regarded as already filled or already understood.  Algebraically, it is the
homology of a quotient chain complex.  Geometrically, it is the natural
language for pairs, CW skeleta, excision, and connecting maps.  This chapter
develops the long exact sequence of a pair and the related exact-sequence
machinery used throughout the rest of the volume.

\section{Relative chains}

Let \(A\subset X\).  The singular chain groups satisfy
\[
C_n(A)\subset C_n(X),
\]
since every singular simplex in \(A\) is also one in \(X\).

\begin{definition}[Relative chain group]
\label{def:viii19-relative-chain}
The relative chain group is
\[
C_n(X,A)=C_n(X)/C_n(A).
\]
\end{definition}

The singular boundary preserves \(C_n(A)\), so it descends to a boundary map
on the quotient.

\begin{proposition}[Relative chain complex]
\label{prop:viii19-relative-complex}
The groups \(C_n(X,A)\) form a chain complex with
\[
\partial[\sigma]=[\partial\sigma].
\]
\end{proposition}

\section{Relative homology}

\begin{definition}[Relative homology]
\label{def:viii19-relative-homology}
\[
H_n(X,A)=H_n(C_*(X,A)).
\]
\end{definition}

A relative cycle may have boundary inside \(A\), because that boundary becomes
zero in the quotient chain complex.

\section{Basic examples}

If \(A=\varnothing\), then
\[
H_n(X,\varnothing)\cong H_n(X).
\]

If \(A=X\), then
\[
H_n(X,X)=0.
\]

If \(X\) is path connected and \(A=\{x_0\}\), then positive-degree relative
homology agrees with reduced homology:
\[
H_n(X,\{x_0\})\cong \widetilde H_n(X)
\qquad(n>0),
\]
with the degree-zero relation controlled by the long exact sequence.

\section{Short exact sequence of chain complexes}

There is a natural short exact sequence
\[
0
\longrightarrow
C_*(A)
\xrightarrow{i_\#}
C_*(X)
\xrightarrow{q_\#}
C_*(X,A)
\longrightarrow
0.
\]

This is exact in every degree.

\section{Connecting homomorphism}

A short exact sequence of chain complexes
\[
0\to A_*\xrightarrow{i}B_*\xrightarrow{p}C_*\to0
\]
produces connecting homomorphisms
\[
\delta_n:H_n(C)\to H_{n-1}(A).
\]

To define \(\delta_n\), take a cycle \(c\in C_n\), lift it to
\(b\in B_n\), observe that \(\partial b\) lies in the image of \(A_{n-1}\),
and take the corresponding homology class in \(H_{n-1}(A)\).

\section{Long exact sequence of a pair}

\begin{theorem}[Long exact sequence of a pair]
\label{thm:viii19-les-pair}
For every pair \(A\subset X\), there is a natural long exact sequence
\[
\cdots
\to H_n(A)
\xrightarrow{i_*}
H_n(X)
\xrightarrow{q_*}
H_n(X,A)
\xrightarrow{\partial}
H_{n-1}(A)
\to\cdots.
\]
\end{theorem}

In low degrees,
\[
H_1(X,A)\xrightarrow{\partial}
H_0(A)\to H_0(X)\to H_0(X,A)\to0.
\]

\section{Exactness meaning}

Exactness at \(H_n(X)\) means
\[
\operatorname{im}i_*=\ker q_*.
\]
Thus an absolute homology class becomes zero relatively exactly when it comes
from the subspace.

Exactness at \(H_n(X,A)\) means
\[
\operatorname{im}q_*=\ker\partial.
\]
Thus a relative class comes from an absolute cycle exactly when its connecting
boundary class vanishes.

\section{Naturality}

A map of pairs
\[
f:(X,A)\to(Y,B),
\qquad
f(A)\subset B,
\]
induces maps on absolute and relative homology.

\begin{theorem}[Naturality of the pair sequence]
\label{thm:viii19-naturality}
The induced maps fit into a commutative diagram of long exact sequences for
\((X,A)\) and \((Y,B)\).
\end{theorem}

\section{The pair \((D^n,S^{n-1})\)}

Because \(D^n\) is contractible and \(S^{n-1}\) has reduced homology concentrated
in degree \(n-1\), the pair sequence gives
\[
H_k(D^n,S^{n-1})
\cong
\begin{cases}
\mathbb Z,&k=n,\\
0,&k\ne n.
\end{cases}
\]

This is the relative calculation underlying cellular chain groups.

\section{Quotient-space interpretation}

For a good pair, such as a CW pair \((X,A)\), collapsing \(A\) to a point
relates relative homology to reduced homology of the quotient.

\begin{theorem}[Relative quotient theorem]
\label{thm:viii19-quotient}
For a CW pair \((X,A)\),
\[
H_n(X,A)
\cong
\widetilde H_n(X/A)
\]
for all \(n\).
\end{theorem}

This is one of the cleanest geometric interpretations of relative homology.

\section{Excision}

\begin{theorem}[Excision]
\label{thm:viii19-excision}
Suppose \(Z\subset A\subset X\) and
\[
\overline Z\subset \operatorname{int}A.
\]
Then inclusion induces an isomorphism
\[
H_n(X\setminus Z,A\setminus Z)
\cong
H_n(X,A).
\]
\end{theorem}

Excision says relative homology ignores material lying completely inside the
interior of the distinguished subspace.

\section{Why excision matters}

Excision allows local decomposition.  It is the technical engine behind:
\begin{itemize}
\item cellular homology,
\item Mayer--Vietoris sequences,
\item local homology,
\item manifold orientation theory.
\end{itemize}

\section{Triples}

Let
\[
B\subset A\subset X.
\]

There is a short exact sequence of relative chain complexes
\[
0\to C_*(A,B)\to C_*(X,B)\to C_*(X,A)\to0.
\]

\begin{theorem}[Long exact sequence of a triple]
\label{thm:viii19-les-triple}
There is a natural long exact sequence
\[
\cdots
\to H_n(A,B)
\to H_n(X,B)
\to H_n(X,A)
\to H_{n-1}(A,B)
\to\cdots.
\]
\end{theorem}

The cellular differential from VIII/18 can be interpreted through such triples
of consecutive skeleta.

\section{Cellular boundary revisited}

For
\[
X^{n-2}\subset X^{n-1}\subset X^n,
\]
the triple sequence contains
\[
H_n(X^n,X^{n-1})
\longrightarrow
H_{n-1}(X^{n-1},X^{n-2}),
\]
which is exactly the cellular boundary map.

Thus the exact-sequence framework explains the construction used in VIII/18.

\section{Relative homology of attached cells}

If
\[
X= A\cup \bigcup_\alpha e^n_\alpha
\]
is obtained from \(A\) by attaching only \(n\)-cells, then
\[
H_k(X,A)
\cong
\begin{cases}
\mathbb Z^{\#\{\text{\(n\)-cells}\}},&k=n,\\
0,&k\ne n.
\end{cases}
\]

This is the cell-by-cell relative calculation behind CW homology.

\section{Relative \(H_0\)}

The group \(H_0(X,A)\) measures components of \(X\) not represented by \(A\).
If \(A\) meets every path component of \(X\), then
\[
H_0(X,A)=0.
\]

\section{Reduced homology from a pair}

For nonempty \(X\) with basepoint \(x_0\),
\[
\widetilde H_n(X)
\cong
H_n(X,\{x_0\})
\]
for \(n>0\), and the pair long exact sequence organizes the degree-zero
adjustment.

\section{Boundary map geometry}

A relative class may be represented by a chain whose boundary lies in \(A\).
The connecting map sends that relative class to the homology class of its
boundary inside \(A\).

This is often the most intuitive way to understand \(\partial\).

\section{Exactness as obstruction theory}

The pair sequence can be read as a succession of obstruction statements:
\begin{itemize}
\item an \(A\)-class dies in \(X\) exactly when it comes from a relative class
one degree higher;
\item a relative class comes from \(X\) exactly when its boundary in \(A\)
vanishes;
\item an \(X\)-class dies relatively exactly when it comes from \(A\).
\end{itemize}

\section{Bridge to homotopy invariance}

Relative homology is functorial, so homotopies of maps of pairs should induce
the same relative maps.  VIII/20 proves the stronger chain-level statement for
singular chains using prism operators and chain homotopies.


\section{Exercises}

\begin{exercise}\label{exr:viii19-01}
Define \(C_n(X,A)\).
\end{exercise}
\begin{hint}
Quotient absolute chains by chains in \(A\).

% VIII-PEDAGOGY-HINT-v1.1 exr:viii19-01 append
The quotient identifies chains differing by a chain supported in \(A\); it does not require every simplex representative to avoid \(A\).
\end{hint}
\begin{solution}
\(C_n(X,A)=C_n(X)/C_n(A)\).
\end{solution}

\begin{exercise}\label{exr:viii19-02}
Why does the boundary descend to the quotient?
\end{exercise}
\begin{hint}
\(C_*(A)\) is a subcomplex.

% VIII-PEDAGOGY-HINT-v1.1 exr:viii19-02 append
Replace \(c\) by \(c+a\), with \(a\in C_n(A)\), and verify that their boundaries differ by a chain in \(C_{n-1}(A)\).
\end{hint}
\begin{solution}
\(\partial C_n(A)\subset C_{n-1}(A)\), so \(\partial[c]=[\partial c]\) is well defined.
\end{solution}

\begin{exercise}\label{exr:viii19-03}
Define \(H_n(X,A)\).
\end{exercise}
\begin{hint}
Take homology of the relative chain complex.

% VIII-PEDAGOGY-HINT-v1.1 exr:viii19-03 append
A relative cycle may have a nonzero absolute boundary, provided that boundary lies in \(C_{n-1}(A)\); apply the quotient differential when testing cycles.
\end{hint}
\begin{solution}
\(H_n(X,A)=H_n(C_*(X,A))\).
\end{solution}

\begin{exercise}\label{exr:viii19-04}
What is \(H_n(X,\varnothing)\)?
\end{exercise}
\begin{hint}
The quotient changes nothing.

% VIII-PEDAGOGY-HINT-v1.1 exr:viii19-04 append
There are no singular simplices with nonempty standard domain mapping to the empty set, so its chain groups vanish.
\end{hint}
\begin{solution}
\(H_n(X)\).
\end{solution}

\begin{exercise}\label{exr:viii19-05}
What is \(H_n(X,X)\)?
\end{exercise}
\begin{hint}
The quotient chain complex is zero.

% VIII-PEDAGOGY-HINT-v1.1 exr:viii19-05 append
Every absolute chain is identified with zero in the quotient by the same chain group; inspect the resulting zero complex in each degree.
\end{hint}
\begin{solution}
Zero.
\end{solution}

\begin{exercise}\label{exr:viii19-06}
Write the chain short exact sequence for a pair.
\end{exercise}
\begin{hint}
Subspace chains, space chains, quotient chains.

% VIII-PEDAGOGY-HINT-v1.1 exr:viii19-06 append
Use the inclusion and quotient maps and verify injectivity, surjectivity, and image equal to kernel degree by degree.
\end{hint}
\begin{solution}
\(0\to C_*(A)\to C_*(X)\to C_*(X,A)\to0\).
\end{solution}

\begin{exercise}\label{exr:viii19-07}
What does this short exact sequence produce?
\end{exercise}
\begin{hint}
Use homological algebra.

% VIII-PEDAGOGY-HINT-v1.1 exr:viii19-07 append
To construct the new arrow, lift a quotient cycle to an absolute chain and take its boundary; that boundary lies in the subspace chain complex.
\end{hint}
\begin{solution}
The long exact sequence of the pair.
\end{solution}

\begin{exercise}\label{exr:viii19-08}
Write the central part of the pair LES.
\end{exercise}
\begin{hint}
Absolute, relative, then one lower subspace group.

% VIII-PEDAGOGY-HINT-v1.1 exr:viii19-08 append
Track degrees on the connecting arrow: it starts at relative degree \(n\) and ends at subspace degree \(n-1\).
\end{hint}
\begin{solution}
\(H_n(A)\to H_n(X)\to H_n(X,A)\xrightarrow{\partial}H_{n-1}(A)\).
\end{solution}

\begin{exercise}\label{exr:viii19-09}
What does exactness at \(H_n(X)\) say?
\end{exercise}
\begin{hint}
Image equals kernel.

% VIII-PEDAGOGY-HINT-v1.1 exr:viii19-09 append
Name the inclusion-induced map \(i_*\) and quotient-induced map \(j_*\); a class is killed by \(j_*\) precisely when it lies in \(\operatorname{im}i_*\).
\end{hint}
\begin{solution}
\(\operatorname{im}H_n(A)\to H_n(X)=\ker(H_n(X)\to H_n(X,A))\).
\end{solution}

\begin{exercise}\label{exr:viii19-10}
What does the connecting map do geometrically?
\end{exercise}
\begin{hint}
Take the boundary of a relative cycle in \(A\).

% VIII-PEDAGOGY-HINT-v1.1 exr:viii19-10 append
Choose \(c\in C_n(X)\) with \(\partial c\in C_{n-1}(A)\); replacing the relative representative must change \(\partial c\) only by an \(A\)-boundary.
\end{hint}
\begin{solution}
It sends a relative class to the homology class of its boundary inside \(A\).
\end{solution}

\begin{exercise}\label{exr:viii19-11}
What is \(H_k(D^n,S^{n-1})\)?
\end{exercise}
\begin{hint}
Use the pair sequence.

% VIII-PEDAGOGY-HINT-v1.1 exr:viii19-11 append
Use \(D^n/S^{n-1}\cong S^n\) for this CW pair, or use the reduced pair sequence to handle the two endpoints when \(n=1\).
\end{hint}
\begin{solution}
\(\mathbb Z\) for \(k=n\), zero otherwise.
\end{solution}

\begin{exercise}\label{exr:viii19-12}
For a CW pair, how is \(H_n(X,A)\) related to \(X/A\)?
\end{exercise}
\begin{hint}
Collapse \(A\).

% VIII-PEDAGOGY-HINT-v1.1 exr:viii19-12 append
Require a nonempty CW subcomplex \(A\) for the usual quotient statement and use reduced homology of the quotient, especially in degree zero.
\end{hint}
\begin{solution}
\(H_n(X,A)\cong\widetilde H_n(X/A)\).
\end{solution}

\begin{exercise}\label{exr:viii19-13}
State excision.
\end{exercise}
\begin{hint}
Remove a subset whose closure lies in the interior of \(A\).

% VIII-PEDAGOGY-HINT-v1.1 exr:viii19-13 append
Write the hypothesis \(\overline Z\subset\operatorname{int}_X(A)\) and identify the inclusion of pairs that induces the claimed isomorphism.
\end{hint}
\begin{solution}
\(H_n(X\setminus Z,A\setminus Z)\cong H_n(X,A)\).
\end{solution}

\begin{exercise}\label{exr:viii19-14}
What is the purpose of excision?
\end{exercise}
\begin{hint}
Ignore interior material of the subspace.

% VIII-PEDAGOGY-HINT-v1.1 exr:viii19-14 append
Apply excision to reduce a cell calculation to \((D^n,S^{n-1})\), checking the interior condition before removing any subset.
\end{hint}
\begin{solution}
It localizes relative calculations and underlies cellular and Mayer--Vietoris arguments.
\end{solution}

\begin{exercise}\label{exr:viii19-15}
Write the inclusions for a triple.
\end{exercise}
\begin{hint}
Three nested spaces.

% VIII-PEDAGOGY-HINT-v1.1 exr:viii19-15 append
Distinguish the three quotient complexes \(C_*(A)/C_*(B)\), \(C_*(X)/C_*(B)\), and \(C_*(X)/C_*(A)\).
\end{hint}
\begin{solution}
\(B\subset A\subset X\).
\end{solution}

\begin{exercise}\label{exr:viii19-16}
What LES comes from a triple?
\end{exercise}
\begin{hint}
Use relative pairs.

% VIII-PEDAGOGY-HINT-v1.1 exr:viii19-16 append
Apply the homology long exact sequence to \(0\to C_*(A,B)\to C_*(X,B)\to C_*(X,A)\to0\), retaining the degree shift.
\end{hint}
\begin{solution}
\(\cdots\to H_n(A,B)\to H_n(X,B)\to H_n(X,A)\to H_{n-1}(A,B)\to\cdots\).
\end{solution}

\begin{exercise}\label{exr:viii19-17}
How does the cellular boundary arise from a triple?
\end{exercise}
\begin{hint}
Use consecutive skeleta.

% VIII-PEDAGOGY-HINT-v1.1 exr:viii19-17 append
Set \((X,A,B)=(X^n,X^{n-1},X^{n-2})\); the triple connecting arrow has target \(H_{n-1}(X^{n-1},X^{n-2})\).
\end{hint}
\begin{solution}
From \(X^{n-2}\subset X^{n-1}\subset X^n\).
\end{solution}

\begin{exercise}\label{exr:viii19-18}
If \(X\) is obtained from \(A\) by \(r\) \(n\)-cells only, what is \(H_n(X,A)\)?
\end{exercise}
\begin{hint}
One generator per cell.

% VIII-PEDAGOGY-HINT-v1.1 exr:viii19-18 append
Collapse the old subcomplex so that the newly attached disks become a wedge of \(n\)-spheres; orient the disks to choose the relative generators.
\end{hint}
\begin{solution}
\(\mathbb Z^r\).
\end{solution}

\begin{exercise}\label{exr:viii19-19}
What is \(H_k(X,A)\) in other degrees in that situation?
\end{exercise}
\begin{hint}
Only \(n\)-cells were added.

% VIII-PEDAGOGY-HINT-v1.1 exr:viii19-19 append
The relative cellular complex has no generators outside degree \(n\); distinguish this relative vanishing from the absolute homology of \(X\).
\end{hint}
\begin{solution}
Zero.
\end{solution}

\begin{exercise}\label{exr:viii19-20}
When is \(H_0(X,A)=0\)?
\end{exercise}
\begin{hint}
Make \(A\) meet every component.

% VIII-PEDAGOGY-HINT-v1.1 exr:viii19-20 append
Examine the cokernel of \(H_0(A)\to H_0(X)\); the inclusion must hit the generator for every path component of \(X\).
\end{hint}
\begin{solution}
When \(A\) meets every path component of \(X\).
\end{solution}

\begin{exercise}\label{exr:viii19-21}
How is reduced homology related to a based pair?
\end{exercise}
\begin{hint}
Use a point subspace.

% VIII-PEDAGOGY-HINT-v1.1 exr:viii19-21 append
Apply the pair sequence with \(A=\{x_0\}\); the map on \(H_0\) is injective, explaining why the degree-one case also works.
\end{hint}
\begin{solution}
For \(n>0\), \(\widetilde H_n(X)\cong H_n(X,\{x_0\})\).
\end{solution}

\begin{exercise}\label{exr:viii19-22}
What maps induce morphisms of pair LESs?
\end{exercise}
\begin{hint}
Maps preserving subspaces.

% VIII-PEDAGOGY-HINT-v1.1 exr:viii19-22 append
Require \(f(A)\subset B\), so the absolute chain map carries the subcomplex being quotiented out into the corresponding target subcomplex.
\end{hint}
\begin{solution}
Maps of pairs \(f:(X,A)\to(Y,B)\).
\end{solution}

\begin{exercise}\label{exr:viii19-23}
Is the pair LES natural?
\end{exercise}
\begin{hint}
Yes; induced maps commute with all arrows.

% VIII-PEDAGOGY-HINT-v1.1 exr:viii19-23 append
For the connecting square, compare \(\partial(f_\#c)\) with \(f_\#(\partial c)\) on a lifted relative cycle.
\end{hint}
\begin{solution}
Yes.
\end{solution}

\begin{exercise}\label{exr:viii19-24}
What is the bridge to VIII/20?
\end{exercise}
\begin{hint}
Homotopies should induce equal homology maps.

% VIII-PEDAGOGY-HINT-v1.1 exr:viii19-24 append
A homotopy of pairs must carry \(A\times I\) into \(B\); this is the condition that allows the prism operator to descend to relative chains.
\end{hint}
\begin{solution}
VIII/20 proves this by a chain homotopy on singular chains.
\end{solution}


\section{Solved problem dossiers}

\begin{problem}[Pair of a disk and boundary]\label{prob:viii19-01}
Compute \(H_*(D^2,S^1)\).
\end{problem}
\begin{solution}
The only nonzero relative group is \(H_2(D^2,S^1)\cong\mathbb Z\).
\end{solution}

\begin{problem}[Interval endpoints]\label{prob:viii19-02}
Compute \(H_1(I,\partial I)\).
\end{problem}
\begin{solution}
Since \(I/\partial I\cong S^1\), the quotient theorem gives \(H_1(I,\partial I)\cong\widetilde H_1(S^1)\cong\mathbb Z\).
\end{solution}

\begin{problem}[Sphere from disk pair]\label{prob:viii19-03}
Use the quotient theorem to rederive \(H_n(D^n,S^{n-1})\).
\end{problem}
\begin{solution}
Collapse \(S^{n-1}\) to a point: \(D^n/S^{n-1}\cong S^n\), so relative homology equals \(\widetilde H_*(S^n)\).
\end{solution}

\begin{problem}[Relative circle to point]\label{prob:viii19-04}
Compute \(H_1(S^1,\{*\})\).
\end{problem}
\begin{solution}
By reduced homology, it is \(\widetilde H_1(S^1)\cong\mathbb Z\).
\end{solution}

\begin{problem}[Relative \(H_0\)]\label{prob:viii19-05}
Let \(X\) have three components and let \(A\) meet exactly two. Compute \(H_0(X,A)\).
\end{problem}
\begin{solution}
One component remains unrepresented by \(A\), so \(H_0(X,A)\cong\mathbb Z\).
\end{solution}

\begin{problem}[Connecting map in disk pair]\label{prob:viii19-06}
What does \(\partial:H_n(D^n,S^{n-1})\to H_{n-1}(S^{n-1})\) do?
\end{problem}
\begin{solution}
Both groups are \(\mathbb Z\), and the connecting map sends the relative fundamental class of the disk to the fundamental class of its boundary sphere; it is an isomorphism up to orientation sign.
\end{solution}

\begin{problem}[Exactness at a sphere]\label{prob:viii19-07}
Why does \(H_n(D^n)=0\) force the connecting map above to be injective?
\end{problem}
\begin{solution}
In the pair LES, exactness gives \(\ker\partial=\operatorname{im}(H_n(D^n)\to H_n(D^n,S^{n-1}))=0\).
\end{solution}

\begin{problem}[Cell attachment pair]\label{prob:viii19-08}
If \(X=A\cup e^n\), compute \(H_*(X,A)\).
\end{problem}
\begin{solution}
The quotient \(X/A\) is \(S^n\), so \(H_n(X,A)\cong\mathbb Z\) and all other relative groups vanish.
\end{solution}

\begin{problem}[Two attached cells]\label{prob:viii19-09}
If two \(n\)-cells are attached to \(A\), compute relative homology.
\end{problem}
\begin{solution}
Collapsing \(A\) gives a wedge \(S^n\vee S^n\) up to the relevant quotient cell structure, so \(H_n(X,A)\cong\mathbb Z^2\).
\end{solution}

\begin{problem}[Excision on a disk]\label{prob:viii19-10}
Remove a small closed disk from the interior of \(D^2\) while keeping an annular subspace inside \(A\). Why can the relative homology remain unchanged?
\end{problem}
\begin{solution}
If the removed set lies with closure inside the interior of the distinguished subspace, excision applies and the inclusion of complements induces an isomorphism.
\end{solution}

\begin{problem}[Triple of skeleta]\label{prob:viii19-11}
For a CW complex, identify the three spaces used to define \(d_n^{\rm cell}\).
\end{problem}
\begin{solution}
Use \(X^{n-2}\subset X^{n-1}\subset X^n\); the triple connecting map runs from \(H_n(X^n,X^{n-1})\) to \(H_{n-1}(X^{n-1},X^{n-2})\).
\end{solution}

\begin{problem}[Absolute class dies relatively]\label{prob:viii19-12}
Suppose \([z]\in H_n(X)\) maps to zero in \(H_n(X,A)\). What does exactness imply?
\end{problem}
\begin{solution}
\([z]\) lies in the image of \(H_n(A)\to H_n(X)\); it is represented, up to homology, by a class from the subspace.
\end{solution}

\begin{problem}[Relative class lifts absolutely]\label{prob:viii19-13}
Suppose \(\xi\in H_n(X,A)\) has \(\partial\xi=0\). What follows?
\end{problem}
\begin{solution}
Exactness says \(\xi\) lies in the image of \(H_n(X)\to H_n(X,A)\).
\end{solution}

\begin{problem}[Subspace class dies in space]\label{prob:viii19-14}
Suppose \(\eta\in H_{n-1}(A)\) maps to zero in \(H_{n-1}(X)\). What follows?
\end{problem}
\begin{solution}
Exactness says \(\eta\) lies in the image of the connecting map \(H_n(X,A)\to H_{n-1}(A)\).
\end{solution}

\begin{problem}[Quotient pair]\label{prob:viii19-15}
For a CW pair \((X,A)\), why is quotient homology reduced?
\end{problem}
\begin{solution}
Collapsing \(A\) creates a distinguished basepoint. Relative chains ignore chains inside \(A\), corresponding to reduced chains of the quotient rather than an extra degree-zero component.
\end{solution}

\begin{problem}[Relative homology of a point pair]\label{prob:viii19-16}
Compute \(H_*(X,X)\).
\end{problem}
\begin{solution}
The quotient chain complex \(C_*(X)/C_*(X)\) is zero, so every relative homology group is zero.
\end{solution}

\begin{problem}[Naturality square]\label{prob:viii19-17}
What commutativity condition must hold for connecting maps under a map of pairs?
\end{problem}
\begin{solution}
If \(f:(X,A)\to(Y,B)\), then \(f_*\partial=\partial f_*\) between \(H_n(X,A)\to H_{n-1}(A)\) and \(H_n(Y,B)\to H_{n-1}(B)\).
\end{solution}

\begin{problem}[Reduced \(H_0\) of two components]\label{prob:viii19-18}
Let \(X\) have two path components and basepoint in one. Compute \(H_0(X,\{x_0\})\).
\end{problem}
\begin{solution}
The pair kills the component containing the basepoint but leaves one independent component, so \(H_0(X,\{x_0\})\cong\mathbb Z\), matching \(\widetilde H_0(X)\).
\end{solution}

\begin{problem}[Relative torus skeleton]\label{prob:viii19-19}
For the torus \(T^2\), what is \(H_2(T^2,T^1)\) in the standard CW structure?
\end{problem}
\begin{solution}
There is one 2-cell, so \(H_2(T^2,T^1)\cong\mathbb Z\).
\end{solution}

\begin{problem}[Relative projective skeleton]\label{prob:viii19-20}
For \(\mathbb RP^n\), compute \(H_n(\mathbb RP^n,\mathbb RP^{n-1})\).
\end{problem}
\begin{solution}
There is one \(n\)-cell, so the relative group is \(\mathbb Z\).
\end{solution}

\begin{problem}[Excision intuition]\label{prob:viii19-21}
Why should relative homology ignore a region lying deep inside \(A\)?
\end{problem}
\begin{solution}
Relative chains already regard chains in \(A\) as zero. Removing material whose closure lies inside \(\operatorname{int}A\) does not affect how chains outside \(A\) meet its boundary.
\end{solution}

\begin{problem}[Boundary as obstruction]\label{prob:viii19-22}
Interpret a nonzero connecting class \(\partial\xi\).
\end{problem}
\begin{solution}
It is the obstruction to representing the relative class \(\xi\) by an absolute cycle in \(X\); exactness says precisely that \(\xi\) lifts absolutely iff this boundary class vanishes.
\end{solution}

\begin{problem}[Long exact sequence source]\label{prob:viii19-23}
What algebraic input produces the LES of a pair?
\end{problem}
\begin{solution}
The short exact sequence of chain complexes \(0\to C_*(A)\to C_*(X)\to C_*(X,A)\to0\).
\end{solution}

\begin{problem}[Bridge to homotopy]\label{prob:viii19-24}
Why should a homotopy of maps of pairs preserve the entire long exact sequence map?
\end{problem}
\begin{solution}
If the two maps induce chain-homotopic chain maps, they induce equal maps on absolute and relative homology, and naturality then preserves all exact-sequence arrows.
\end{solution}
